\begingroup
\raggedbottom
\subsection{What the results show}
\label{sec:benchmark-interpretation}

The results distinguish two ways of using the same reference:
\begin{itemize}
  \item \textbf{Index:} build the reference index first, restart PostgreSQL,
  then measure the query, including access to the stored index.
  \item \textbf{No index:} measure construction of a temporary reference index
  and the query together. The query batch shares that temporary index.
\end{itemize}

\Needspace{13\baselineskip}
\subsubsection{What an indexed minimizer comparison measures}
\label{sec:bench-indexed-minimizers}

Scalar and SIMD-ref have similar indexed times in this example, where they
reuse the same kind of reference index and share most of the query pipeline.

\BenchmarkSlice{aws-2-k15-w10}
  {mpf-scalar-scalar-block,mpf-simdref-scalar-block}
  {index-2-k15-w10-dna4-m7a,index-2-k15-w10-dna4-m7i}
  {Two minimizer implementations with MPF, Scalar-DP and Block held fixed.}
  {tab:bench-indexed-minimizers}

\warn{An indexed row does not measure reference-index construction}{
The benchmark requires a compatible persisted \emph{reference} index. It builds
that index before restarting PostgreSQL, and does not rebuild its minimizers
inside the measured query. The query sequences themselves have no persisted
minimizer index: their minimizers are extracted to search the reference.
Changing the minimizer therefore does not measure reference construction
again, but it need not produce an identical query execution.}

The distinction follows the reference-read policy and query-side preparation
described in Section~\ref{sec:index-read-policy}.
It is also checked by the campaign: the query fixture has no generated indexes,
and accepted indexed records report query-sketch work alongside zero target-index
recomputation. Scalar and SIMD-ref can share a persisted hash representation,
whereas bounded ntHash can change both the query sketch and the reference
lookup representation (Sections~\ref{sec:configured-cpu-minimizers}
and~\ref{sec:packed-minimizer-representation}).
Close timings are therefore consistent with substantial shared work, rather
than proof that the complete execution was measured twice.

In each small table, the upper duration is A, the complete query measurement,
and the lower duration is B, the residual after removing the recorded CIGAR/MAPQ
intervals. Cells, comparison symbols and colours are copied from the full table.
The detailed timing definition and remaining markers are explained in
Section~\ref{sec:benchmark-results}.

\Needspace{13\baselineskip}
\subsubsection{Paying for the reference once}
\label{sec:bench-index-benefit}

The preferred BioseqDB configuration beats minimap2 when the reference index
is available, but loses when temporary construction is included in this example.

\BenchmarkSlice{aws-2-k15-w10}
  {minimap2,mpf-bounded-simd-block}
  {index-2-k15-w10-dna4-m7i,noindex-2-k15-w10-dna4-m7i}
  {The preferred configuration on M7I, with and without a prepared reference index.}
  {tab:bench-index-benefit}

A persistent index moves reference minimizer extraction and organization out of
the query. BioseqDB can then read its prepared lookup structures, as shown in
Figures~\ref{fig:index-join-read-paths} and~\ref{fig:index-control-data-plane}.
Without that index it must construct a transient representation before the batch
can use it. This is consistent with the database use case in
Section~\ref{sec:index-storage-role}: preparation can be worthwhile when the
same reference serves many later queries.

The difference between these columns is not an isolated measurement of building
the persistent index. Temporary construction and persisted lookup use different
paths. Durable publication, described in Section~\ref{sec:index-file-consistency},
is outside both query measurements. The result supports reuse of a prepared
reference, but does not by itself give its construction cost or the number
of queries needed to recover that cost.

\Needspace{13\baselineskip}
\subsubsection{A smaller reference can have a more expensive lookup path}
\label{sec:bench-wide-hash}

The scalar MPF configuration takes much longer on the smaller reference with
longer k-mers than on the larger reference with shorter k-mers.

\BenchmarkSlice{aws-4-k15-w10}
  {mpf-scalar-scalar-block,mpf-bounded-simd-block}
  {index-4-k15-w10-dna4-m7a,index-4-k15-w10-dna4-m7i}
  {Indexed DNA4 with the shorter k-mer profile.}
  {tab:bench-short-hash}
\BenchmarkSlice{aws-2-k20-w15}
  {mpf-scalar-scalar-block,mpf-bounded-simd-block}
  {index-2-k20-w15-dna4-m7a,index-2-k20-w15-dna4-m7i}
  {Indexed DNA4 with the longer k-mer profile, despite a smaller reference.}
  {tab:bench-wide-hash}

The current implementation offers a concrete explanation through the hash
representation and its reader. The integer-hash profile for shorter k-mers fits
the compact direct lookup described in Section~\ref{sec:index-direct-lookup}.
The longer integer hashes require the wider representation and cannot use that
compact reader. The alternative reader scans stored metadata to prepare and
emit anchors. Successive groups of queries can revisit this metadata, as explained in
Section~\ref{sec:index-control-scan-waves}. Mapping an index file therefore does
not make every lookup equally cheap.

The bounded configuration is a useful counterexample: its compact hash profile
retains eligibility for direct lookup, and the longer-k-mer case remains fast here.
The bounded sketcher has a specialised default-profile kernel and a runtime
parameter path (Section~\ref{sec:configured-cpu-minimizers}), but a generic
sketch loop alone does not explain a large delay in a query that already has
its reference index. These measurements also span campaign builds, so
current source identifies a credible mechanism rather than a measured breakdown
of every historical run.

A larger possible k-mer vocabulary does not mean an array of every possible
k-mer is allocated. Index size depends on the observed minimizers, their stored
width and their occurrence lists. Increasing the window can reduce minimizer
density, while changing k also changes matches and repetition
(Section~\ref{sec:minimizer-definition} and Equation~\ref{eq:minimizer-anchor-count}).
Longer k-mers therefore do not imply either a larger index or swapping in every case.

\Needspace{23\baselineskip}
\subsubsection{SIMD and RMQ are workload-dependent choices}
\label{sec:bench-chaining-choice}

The preferred SIMD configuration is competitive, but the fastest chaining
choice changes with the workload.

\BenchmarkSlice[columns]{aws-2-k15-w10}
  {mpf-bounded-scalar-block,mpf-bounded-simd-block,mpf-bounded-rmq-block}
  {index-2-k15-w10-dna4-m7a,index-2-k15-w10-dna4-m7i}
  {Chaining choices with MPF, bounded minimizers and Block fixed.}
  {tab:bench-chaining-small}

Here the complete-query times are close. Scalar chaining already avoids much
unproductive work through bounded predecessor searches and stopping rules
(Algorithm~\ref{alg:chaining-b1}). SIMD must also arrange coordinates and scores
for vector processing, and its backtracking needs additional storage
(Section~\ref{sec:chaining-simd-forward} and
Algorithm~\ref{alg:chaining-simd-backtrack}). Those costs can offset a cheaper
score calculation. Shared index access and anchor generation can also hide a
change in chaining time. The totals do not establish which effect dominates,
or prove that compiler autovectorization explains a scalar win.

\BenchmarkSlice[columns]{aws-4-k15-w10}
  {mpf-bounded-scalar-block,mpf-bounded-simd-block,mpf-bounded-rmq-block}
  {index-4-k15-w10-dna4-m7a,index-4-k15-w10-dna4-m7i}
  {The same choices on the larger DNA4 reference.}
  {tab:bench-chaining-large}

The larger case favours Scalar-DP. This is also a comparison of stopping
policies: the SQL SIMD-DP choice uses bounded approximate batches, rather than
the exact scalar predecessor traversal (Algorithm~\ref{alg:chaining-b5}).
RMQ replaces predecessor scanning with tree maintenance and range queries.
Its benefit depends on how much scanning it saves. Its current candidate
selection also differs from minimap2 RMQ
(Section~\ref{sec:chaining-bioseqdb-rmq-summary} and
Table~\ref{tab:chaining-bioseqdb-rmq-differences}). These alternatives need not
perform identical work or produce identical intermediate regions.

\BenchmarkSlice[columns]{local}
  {mpf-bounded-scalar-block,mpf-bounded-simd-block,mpf-bounded-rmq-block}
  {index-2-k20-w15-dna4-local}
  {A local case in which SIMD-DP wins and RMQ loses.}
  {tab:bench-chaining-local}

The local result prevents a general conclusion that SIMD helps less on the
Apple machine than on AWS. For matched, complete five-repeat comparisons with
MPF, bounded minimizers, Block and DNA4 fixed, the median AWS differences are
small. Restricting the comparison to measurements from the same source file
preserves that conclusion. The local coverage is too sparse for an equivalent
aggregate. These slices support selecting a chaining method using matched
workloads.

\Needspace{13\baselineskip}
\subsubsection{Vector width alone does not explain the CPU results}
\label{sec:bench-vector-width}

Scalar and reference SIMD minimizers can remain close on both the local
machine and the AWS machines after reference construction is removed.

\BenchmarkSlice[columns]{local}
  {mps-scalar-simd-block,mps-simdref-simd-block}
  {index-2-k15-w10-dna4-local}
  {Local indexed minimizer choices with MPS, SIMD-DP and Block fixed. Both are single samples.}
  {tab:bench-vector-width}

The local cells are close single observations, not a repeatability estimate.
An indexed query extracts minimizers from the queries rather than sketching
the entire reference again. A faster sketcher can therefore have little effect
on complete-query time. Moreover, the reference SIMD implementation has shared
rescan work when lanes lose their minima. Vector width does not remove that
cost (Section~\ref{sec:rewrite-naive-simd}). The bounded implementation uses
different work organization, described in
Section~\ref{sec:configured-cpu-minimizers}.

The current bounded kernel also interleaves two streams on the four-lane path
and one on the eight-lane path, keeping the number of logical streams the same
(Section~\ref{sec:minimizer-width-and-query-work}). This changes instruction
scheduling and register demand. It does not make the two processors equally
fast. The earlier stream experiments in Section~\ref{sec:v11-v12} explain
the motivation. They do not establish a universal SIMD speedup across the
different CPUs, index readers and memory conditions in this campaign.

\Needspace{13\baselineskip}
\subsubsection{Ambiguous bases expose work beyond minimizer extraction}
\label{sec:bench-dna5}

The local DNA5 slowdown is much larger in the complete no-index measurement
than in its residual measurement.

\BenchmarkSlice{local}
  {mpf-bounded-simd-block}
  {index-2-k15-w10-dna4-local,index-2-k15-w10-dna5-local,noindex-2-k15-w10-dna4-local,noindex-2-k15-w10-dna5-local}
  {The preferred local configuration on DNA4 and DNA5.}
  {tab:bench-dna5}

Ambiguous bases change the workload at several stages. The bounded DNA5
sketcher includes N in the k-mers it hashes. It does not simply skip every k-mer
containing N like the minimap2-style ambiguity rule. Its hashes and selected
anchors can therefore change along with the alphabet
(Section~\ref{sec:configured-cpu-minimizers}). The CPU kernel supports this
profile, so these rows are not evidence of an unnoticed scalar replacement.
DNA5 also uses a different packed representation
(Section~\ref{sec:sequence-bit-layout}).

However, the large gap that disappears from B points towards time covered by
the alignment/CIGAR intervals, not solely towards an expensive N test in the
sketcher. The fused alignment work described in
Section~\ref{sec:final-alignment-fused-work} can change when anchors or exact
matching runs change. The totals do not isolate that mechanism from
other alignment work. They justify investigating the downstream alignment
path before attributing the whole slowdown to DNA5 minimizer extraction.

\Needspace{13\baselineskip}
\subsubsection{A fast residual does not mean a fast complete query}
\label{sec:bench-a-b}

This MPS configuration loses to minimap2 on complete time even though its
residual time is lower.

\BenchmarkSlice{aws-2-k20-w15}
  {minimap2,mps-scalar-simd-block}
  {index-2-k20-w15-dna4-m7a,index-2-k20-w15-dna4-m7i}
  {An indexed comparison whose direction changes between A and B.}
  {tab:bench-a-b}

The CIGAR interval includes base dynamic programming, trace storage, traceback
and result construction, not just conversion of an existing alignment into
text (Section~\ref{sec:final-alignment-fused-work}). A large A/B separation
therefore identifies substantial wall time covered by that instrumented work.
It is not evidence of a costly output copy alone.

There is a further limit under parallel execution: one query can be aligning
while another is still seeding or chaining
(Section~\ref{sec:arch-work-distribution}). Removing the union of alignment
intervals also removes wall time during which those other tasks ran. B is
neither an isolated sum of early-stage costs nor a prediction of running the
query with CIGAR and MAPQ disabled. Close A and B show little wall time covered
by the removed intervals. They do not identify one remaining stage as the bottleneck.

The historical divergence-sort problem in
Section~\ref{sec:minimizer-divergence-history} is a different issue.
B removes the recorded CIGAR and MAPQ intervals. It does not subtract a separate
timer for divergence estimation, which precedes final alignment for each query.
Divergence work can nevertheless overlap another query's alignment interval.
The current estimator uses prefix counts
(Section~\ref{sec:divergence-prefix-counts}). The reversal between A and B does
not demonstrate a recurrence of the earlier sorting problem.

\Needspace{13\baselineskip}
\subsubsection{Block Aligner improves complete query time}
\label{sec:bench-base-aligner}

Block Aligner improves complete time in this matched comparison, while the
residual changes are smaller.

\BenchmarkSlice{aws-2-k15-w10}
  {mpf-bounded-simd-block,mpf-bounded-simd-ksw2}
  {index-2-k15-w10-dna4-m7a,index-2-k15-w10-dna4-m7i,index-2-k15-w10-dna4-m4}
  {Changing only the requested base aligner in the preferred configuration.}
  {tab:bench-base-aligner}

This pattern is consistent with a benefit in the fused alignment work.
Block Aligner follows an adaptive block through the scoring matrix, whereas
KSW2 uses its banded dynamic programming and stopping rules. Their gap models
and extension cutoffs differ
(Section~\ref{sec:final-alignment-cpu-contracts}). Fewer computed cells can
matter as much as the speed of one cell update.

Matched five-repeat DNA4 comparisons favour Block on the shorter-k-mer AWS
workload, but occasional reversals remain elsewhere. A reversal accompanied
by a large change in B is weak evidence about the base aligner alone: residual
wall time can change with concurrent overlap and execution conditions as well
as other work. The current pipeline uses a small
adaptive block range, described in
Section~\ref{sec:final-alignment-cpu-contracts}. Changing that range may trade
more score cells against less block-management work, but read length alone
does not determine the best setting: the size and divergence of each region
being extended matter. These tables do not contain a block-size sweep.

\Needspace{13\baselineskip}
\subsubsection{MPF and MPS change both lookup and downstream work}
\label{sec:bench-index-layout}

MPF is faster in this matched example, but the no-index comparison does not
measure persistent-index construction in isolation.

\BenchmarkSlice{aws-2-k15-w10}
  {mpf-bounded-simd-block,mps-bounded-simd-block}
  {index-2-k15-w10-dna4-m7a,noindex-2-k15-w10-dna4-m7a}
  {MPF and MPS with the same minimizer, chaining and alignment choices.}
  {tab:bench-index-layout}

MPF retains fragment-oriented lookup data. MPS groups occurrences over the
sequence (Figure~\ref{fig:index-mpf-mps}). Their readers expose different
metadata, occurrence grouping and access patterns. The hash-major MPS reader
can read selected streams rather than every stored position
(Section~\ref{sec:index-mps-selective-reads}), but that advantage does not
guarantee a cheaper complete query than MPF direct lookup. Matched five-repeat
bounded DNA4 comparisons usually favour MPF for indexed execution. The
no-index differences are smaller and are sometimes close.

Persistent MPS construction includes sorting and may spill runs to temporary
files (Figure~\ref{fig:index-mps-merge}). That explains a possible preparation
cost, not a cost automatically charged to an indexed query. The sorter already
keeps runs in memory when its memory budget allows and spills when needed.
Temporary construction,
persistent publication and later reading are different operations
(Sections~\ref{sec:index-construction} and~\ref{sec:index-file-consistency}).
An index/no-index subtraction cannot rank their durable build times.

\BenchmarkSlice{aws-4-k20-w15}
  {mpf-bounded-simd-block,mps-bounded-simd-block}
  {index-4-k20-w15-dna5-m7a,index-4-k20-w15-dna5-m7i}
  {The large DNA5 MPS slowdown appears on several machine classes.}
  {tab:bench-mps-dna5}

The MPS anomaly is not unique to M7A. Much of its additional elapsed time is
covered by the intervals removed from B, so persistent spill sorting does not
explain the gap. Current DNA5 MPS uses the grouped reader and a serial query
loop with fragmented target access, whereas MPF can use a shared resident
target and parallel batch (Section~\ref{sec:index-reader-execution}). Alignment
requests through the PostgreSQL view copy ranks through cached fragment access
and per-symbol extraction (Section~\ref{sec:sequence-pg-rank-copy}). This is
different work from reading a shared resident target, and provides a possible
connection between the index route and the fused alignment interval.
The packed fragment slots are separate from the decompressed stripe cache,
which was disabled in the original benchmark campaign
(Section~\ref{sec:pg-storage-caches}).

Each MPS cell comes from one measured batch in the FAST fallback procedure,
whereas the MPF cells have five repeats. The current source provides a plausible
explanation, but there is no stage trace from each historical binary.
The table cannot assign the whole gap
to fragmentation, serial execution or memory pressure separately.

\Needspace{13\baselineskip}
\subsubsection{Machine comparisons include memory and I/O pressure}
\label{sec:bench-machine-pressure}

The minimap2 baseline itself slows sharply on C7G, so its large BioseqDB
comparison factors also reflect a stressed reference run.

\BenchmarkSlice{aws-4-k20-w15}
  {minimap2,mpf-bounded-simd-block}
  {index-4-k20-w15-dna4-m7a,index-4-k20-w15-dna4-m7i,index-4-k20-w15-dna4-m4,index-4-k20-w15-dna4-c7g}
  {The same indexed workload across the AWS machine classes.}
  {tab:bench-machine-pressure}

The C7G minimap2 record reports severe paging during measurement.
Its available physical memory was nearly exhausted, and a large part of the
helper's address space was swapped out. C7G also has less physical RAM than
the other AWS classes in this campaign. The residency distinctions in
Section~\ref{sec:index-residency} matter directly: accessing a stored index
can wait for data even when no index is rebuilt.

M4 is slower in many comparisons, but is not uniformly the slowest machine.
Its older CPU and lower instance-level EBS bandwidth can both matter when
working data must be fetched from disk. The common provisioned volume settings
do not make the effective I/O paths equal.\autocite{aws_ebs_limits}
Neither the inventory nor the elapsed times measure a universal swap penalty
for each processor. The C7G baseline measures execution under severe memory
pressure. Its large comparison factor should therefore not be presented as a
BioseqDB speedup that applies independently of the machine and its memory
conditions.

\Needspace{13\baselineskip}
\subsubsection{An indexed query can lose to temporary construction}
\label{sec:bench-index-reversal}

The C7G measurement is slower with a persisted index than without one.

\BenchmarkSlice{aws-2-k20-w15}
  {minimap2,mpf-simdref-simd-block}
  {index-2-k20-w15-dna4-c7g,noindex-2-k20-w15-dna4-c7g}
  {The apparently reversed index benefit for MPF and SIMD-ref.}
  {tab:bench-index-reversal}

The indexed record explicitly reports persisted MPF decoding and query-sketch
jobs with no reference-index build or recomputation. Restarting PostgreSQL
removes process-resident state. It does not remove the persisted index. The
slow time must therefore be interpreted as querying and accessing that index,
rather than silently constructing it again.

The two modes need not perform identical lookup work. A temporary index can
use an in-memory representation, whereas a wide-hash persisted reader can
repeatedly traverse stored controls
(Section~\ref{sec:index-control-scan-waves}). Memory residency adds another
possible difference (Section~\ref{sec:index-residency}). These cells also come
from different campaign replicas with different execution histories. The
generator selects the lowest approved A independently for each cell. The
table therefore shows a real recorded reversal, but does not establish a timing
bug or show that removing the index repeatedly improves performance on one
fixed deployment.

\Needspace{13\baselineskip}
\subsubsection{The longest runs need a resource explanation before a kernel ranking}
\label{sec:bench-long-runs}

Several wide-hash MPF observations remain extremely slow even after the
alignment intervals are removed.

\BenchmarkSlice[columns]{aws-4-k20-w15}
  {mpf-scalar-scalar-block,mpf-simdref-rmq-block}
  {index-4-k20-w15-dna4-c7g}
  {Examples of the longest indexed observations. Superscript s denotes a single sample.}
  {tab:bench-long-runs}

Each of these very long measurements comes from a single FAST observation.
Their close A and B place most of the delay outside the removed alignment intervals,
but do not separate index traversal, anchor work, page faults and scheduling.
The repeated control scans described in Section~\ref{sec:index-control-scan-waves}
can amplify slow memory or storage access: the same large metadata region is
visited for successive query waves. Changing the final base aligner cannot
remove that preceding work.

The campaign logs also contain a backend that appeared idle while its disk-read
counter continued to advance and the kernel reported a page wait. That example
occurred during untimed preparation. It explains why low CPU use alone is not
evidence of a deadlock, but cannot diagnose these timed cells. The admission
path in Section~\ref{sec:arch-admission} and
Algorithm~\ref{alg:arch-wave-scheduling} returns a decision rather than
sleeping until memory becomes available. A denied reservation makes the caller
reduce, serialize or reject work. This rules out treating the admission API
itself as a demonstrated indefinite memory wait.

The available records contain no per-repeat trace proving a TBB deadlock in
these completed runs, nor a uniform paging trace proving that every long cell had the same
cause. A controlled repeat would need the actual reader, wave progress,
page-fault/I/O counters and worker stacks together. Until then, these observations
reveal a serious robustness limit of the measured path. They should not
determine a general Scalar-DP versus RMQ ranking.

\Needspace{13\baselineskip}
\subsubsection{A larger GPU does not remove the host part of the pipeline}
\label{sec:bench-gpu-performance}

The local Vulkan route remains slower than minimap2 when the reference index
is already available.

\BenchmarkSlice{local}
  {minimap2,mpf-vulkan-vulkan-block}
  {index-2-k15-w10-dna4-local,index-2-k20-w15-dna4-local}
  {The local integrated GPU with Vulkan seeding and chaining.}
  {tab:bench-gpu-local}

The local benchmark used unified CPU/GPU memory, so its result cannot be
explained simply as using CPU RAM instead of separate GPU VRAM. The published
route includes host work, GPU submissions and synchronization. Its selected
base aligner runs on the CPU, and the longer-k-mer local record also identifies
CPU backtracking (Section~\ref{sec:chaining-gpu-campaign-boundary}). The staged minimizer
schedule in Algorithm~\ref{alg:vulkan-sketch-host-schedule} illustrates why
fast device computation need not imply fast complete-query execution.

\BenchmarkSlice{vastai}
  {minimap2,mps-vulkan-vulkan-block}
  {index-2-k15-w10-dna4-vastai,index-2-k20-w15-dna4-vastai}
  {The MPS Vulkan route and the same-host minimap2 baseline on an RTX 5090.}
  {tab:bench-gpu-vast}

The discrete-GPU results are more mixed: the longer-k-mer MPS case approaches
the baseline on complete time, while the other case remains much slower.
These are single observations. They use a different host and software build
from the local campaign, so they do not isolate the effect of replacing the GPU.
The Vast.ai build places GPU working buffers in device-local memory but retains
CPU-cached arrays for host consumers. The campaign diagnostics found that
moving host-read chaining data towards the GPU could itself slow backtracking.
The useful design rule is to place each array near the work that consumes it
and reduce transfers between stages.

G3SA targets long-read alignment too. Its results do not imply that GPUs are
useful only for short reads. Its released pipeline uses multiple kernels and
retains major intermediate arrays on the device, rather than placing the
whole mapper in one kernel.\autocite{g3sa_paper,g3sa_implementation}
Algorithm~\ref{alg:g3sa-sketch-kernel} describes one such stage. This suggests
a direction for reducing host/device interaction, but its different pipeline
does not predict the speedup of the hybrid PostgreSQL implementation measured here.

\Needspace{13\baselineskip}
\subsubsection{GPU capacity failures have more than one cause}
\label{sec:bench-gpu-capacity}

The larger MPS GPU cases fail even though the MPF cases complete on the same device.

\BenchmarkSlice{vastai}
  {mpf-vulkan-vulkan-block,mps-vulkan-vulkan-block}
  {index-4-k15-w10-dna4-vastai,index-4-k20-w15-dna4-vastai}
  {Completed MPF observations and MPS preparation failures on the RTX 5090.}
  {tab:bench-gpu-capacity}

The diagnosed MPS failure is a storage-buffer binding limit. Its combined
hash/offset/count control array exceeds the device's allowed storage-buffer
range, although the position buffer fits. These are index controls of the
kind separated from position streams in Figure~\ref{fig:format6-streams},
not just one raw sequence fragment. Available total GPU memory does not
establish that one binding can represent the whole array. Splitting the work
into bounded buffer windows is required even on a device with more VRAM.
The table does not contain a successful full-workload measurement of that repair.

For the longer-k-mer profile, the detailed buffer diagnostics come from a
companion KSW2 run using the same MPS index profile. The displayed Block run
records a more general admission failure.

Other absent cells have different meanings. The measured Vast.ai build rejected
the DNA5 sequence type during setup. Missing or failed execution evidence
does not constitute a slow timing. These distinctions follow the admission
principle in Section~\ref{sec:arch-route-admission}: a selected route must be
supported and proved to have run before its elapsed time is compared.

\Needspace{13\baselineskip}
\subsubsection{Search dominates this BLAST measurement}
\label{sec:bench-external}

BLAST remains much slower than the mapping-oriented tools even when its native
reference database is prepared beforehand.

\BenchmarkSlice[columns]{external-2-k15-w10}
  {minimap2,bioseq,winnowmap,lexicmap,blastn}
  {index-2-k15-w10-dna4-m7a,noindex-2-k15-w10-dna4-m7a}
  {The preferred BioseqDB tuple and the external aligners on M7A.}
  {tab:bench-external}

The external timing includes launching the process and producing its output,
but the BLAST stage evidence places most elapsed time inside word
finding and preliminary search. The long indexed time is therefore not well
explained as executable startup alone. BLAST's seed search and extension
strategy differs from the minimizer/chain pipeline introduced in
Section~\ref{sec:seed-chain-extend}. Preparing a native database does not
remove that query work. Increasing the word length can reduce it, so the
longer-k-mer profile is not universally slower across tools.

Winnowmap also benefits strongly from a prepared index in this slice, while
LexicMap's indexed search remains slower than minimap2 on this workload.
The comparison rules in Section~\ref{sec:external-aligner-results} remain
important: native indexing, reference partitioning and ambiguity handling
differ between tools. LexicMap has no equivalent minimizer-window parameter,
and the B intervals do not all cover the
same fused work as BioseqDB. The useful result is the complete time for the
stated workload and validation rules. It does not establish interchangeable
alignment semantics or an isolated kernel speed ranking.
\clearpage
\endgroup
